%%
%% Test: Persian Locale (fa-IR) — v1.1
%% Purpose: Verify Persian localization of LaTeX names and cleveref
%% 

\documentclass{matinbook}

\title{تست بومی‌سازی فارسی — نسخه ۱.۱}
\author{تیم MatinBook}
\date{۱۴۰۵}

\begin{document}

\maketitle

\tableofcontents

\chapter{بررسی بومی‌سازی فارسی}
\label{chap:test-locale}

این سند، بومی‌سازی فارسی \texttt{fa-IR} را بررسی می‌کند.

\section{نام‌های ساختاری}
\label{sec:structure-names}

\begin{itemize}
    \item \verb|\contentsname|: \contentsname
    \item \verb|\listfigurename|: \listfigurename
    \item \verb|\listtablename|: \listtablename
    \item \verb|\bibname|: \bibname
    \item \verb|\indexname|: \indexname
    \item \verb|\chaptername|: \chaptername
    \item \verb|\appendixname|: \appendixname
    \item \verb|\proofname|: \proofname
    \item \verb|\abstractname|: \abstractname
    \item \verb|\figurename|: \figurename
    \item \verb|\tablename|: \tablename
\end{itemize}

\section{ارجاعات cleveref}
\label{sec:cleveref-names}

\begin{equation}
    \label{eq:locale-test}
    f(x) = x^2
\end{equation}

\begin{figure}[h]
\centering
\rule{3cm}{2cm}
\caption{شکل نمونه}
\label{fig:locale-test}
\end{figure}

\begin{table}[h]
\centering
\caption{جدول نمونه}
\label{tab:locale-test}
\begin{tabular}{cc}
\toprule
الف & ب \\
\midrule
۱ & ۲ \\
\bottomrule
\end{tabular}
\end{table}

\begin{theorem}[قضیه نمونه]
    برای هر عدد حقیقی $x$، $x^2 \geq 0$.
\end{theorem}
\label{thm:locale-test}

ارجاعات:

\begin{itemize}
    \item معادله: \cref{eq:locale-test}
    \item شکل: \cref{fig:locale-test}
    \item جدول: \cref{tab:locale-test}
    \item قضیه: \cref{thm:locale-test}
    \item فصل: \cref{chap:test-locale}
    \item بخش: \cref{sec:structure-names}
\end{itemize}

\section{قضیه و اثبات}
\label{sec:theorem-proof}

\begin{theorem}[قضیه نمونه]
    برای هر عدد حقیقی $x$، داریم $x^2 \geq 0$.
\end{theorem}

\begin{proof}
    چون مربع هر عدد حقیقی نامنفی است، نتیجه حاصل می‌شود.
\end{proof}

\section{نتیجه‌گیری}
\label{sec:conclusion}

اگر این سند بدون خطا کامپایل شود:

\begin{itemize}
    \item ✅ نام‌های ساختاری فارسی هستند
    \item ✅ ارجاعات cleveref فارسی هستند
    \item ✅ محیط‌های قضیه و اثبات کار می‌کنند
\end{itemize}

\textbf{وضعیت تست:} قبول

\end{document}
